\begin{blocksection}
\question \textbf{Labels at a depth.} Implement \lstinline{labels_at_depth},
which takes a tree \lstinline{t} with integer labels and a nonnegative integer
\lstinline{k}. Return a new Python list containing the labels exactly
\lstinline{k} edges below the root, in left-to-right branch order. The root is
at depth 0. If no nodes exist at that depth, return an empty list.

Use recursion and \lstinline{extend} to add each recursive result to the local
\lstinline{result} list. Do not change any tree labels or branches.

\begin{lstlisting}
def labels_at_depth(t: Tree, k: int) -> list[int]:
    """
    >>> t = Tree(1, [Tree(2, [Tree(4), Tree(5)]),
    ...              Tree(3, [Tree(6)])])
    >>> labels_at_depth(t, 0)
    [1]
    >>> labels_at_depth(t, 1)
    [2, 3]
    >>> labels_at_depth(t, 2)
    [4, 5, 6]
    >>> labels_at_depth(t, 3)
    []
    >>> labels_at_depth(Tree(9), 1)
    []
    """
    if _________________________________:

        return _________________________

    result = []

    for b in ___________________________:

        result.extend(__________________)

    return result
\end{lstlisting}
\end{blocksection}

\begin{solution}
\begin{lstlisting}
def labels_at_depth(t: Tree, k: int) -> list[int]:
    if k == 0:
        return [t.label]
    result = []
    for b in t.branches:
        result.extend(labels_at_depth(b, k - 1))
    return result
\end{lstlisting}
At depth 0, return a one-element list containing the root's label. Otherwise,
the requested nodes are at depth \lstinline{k - 1} within each branch.
Each call creates its own result list. \lstinline{extend} mutates that list,
but the input tree is only read. At a leaf with positive \lstinline{k}, the
loop has no iterations, so the result is empty.
\end{solution}

\begin{blocksection}
Explain what would change if \lstinline{extend} were replaced by
\lstinline{append}. Why would \lstinline{return result.extend(...)} be wrong?
\begin{solution}[.65in]
\lstinline{append} adds each recursive result as a single list element,
introducing nested lists. For example, at depth 1 the result would be
\lstinline{[[2], [3]]}. \lstinline{extend} adds the elements individually.
It returns \lstinline{None}, so we must return \lstinline{result} separately.
\end{solution}
\end{blocksection}

\begin{questionmeta}
Suggested Time: 10 min; Difficulty: Medium
\begin{itemize}
    \item Ask what the recursive call returns before filling in the loop.
    Decreasing \lstinline{k} measures the remaining distance, not a label value.
    \item Draw a fresh result list for each recursive call. Mutation occurs in
    those lists; no assignment targets a tree attribute or its branches list.
    \item Trace an uneven tree or a depth beyond the leaves to explain why a
    separate leaf base case is unnecessary.
\end{itemize}
\end{questionmeta}
